% Options for packages loaded elsewhere
\PassOptionsToPackage{unicode}{hyperref}
\PassOptionsToPackage{hyphens}{url}
\PassOptionsToPackage{dvipsnames,svgnames,x11names}{xcolor}
\documentclass[
  11pt,
]{article}
\usepackage{xcolor}
\usepackage[a4paper,margin=19mm]{geometry}
\usepackage{amsmath,amssymb}
\setcounter{secnumdepth}{-\maxdimen} % remove section numbering
\usepackage{iftex}
\ifPDFTeX
  \usepackage[T1]{fontenc}
  \usepackage[utf8]{inputenc}
  \usepackage{textcomp} % provide euro and other symbols
\else % if luatex or xetex
  \usepackage{unicode-math} % this also loads fontspec
  \defaultfontfeatures{Scale=MatchLowercase}
  \defaultfontfeatures[\rmfamily]{Ligatures=TeX,Scale=1}
\fi
\usepackage{lmodern}
\ifPDFTeX\else
  % xetex/luatex font selection
  \setmainfont[]{Be Vietnam Pro}
  \setmonofont[]{Menlo}
\fi
% Use upquote if available, for straight quotes in verbatim environments
\IfFileExists{upquote.sty}{\usepackage{upquote}}{}
\IfFileExists{microtype.sty}{% use microtype if available
  \usepackage[]{microtype}
  \UseMicrotypeSet[protrusion]{basicmath} % disable protrusion for tt fonts
}{}
\makeatletter
\@ifundefined{KOMAClassName}{% if non-KOMA class
  \IfFileExists{parskip.sty}{%
    \usepackage{parskip}
  }{% else
    \setlength{\parindent}{0pt}
    \setlength{\parskip}{6pt plus 2pt minus 1pt}}
}{% if KOMA class
  \KOMAoptions{parskip=half}}
\makeatother
\usepackage{longtable,booktabs,array}
\usepackage{caption}
\captionsetup[table]{skip=6pt}
\newcounter{none} % for unnumbered tables
\usepackage{calc} % for calculating minipage widths
% Correct order of tables after \paragraph or \subparagraph
\usepackage{etoolbox}
\makeatletter
\patchcmd\longtable{\par}{\if@noskipsec\mbox{}\fi\par}{}{}
\makeatother
% Allow footnotes in longtable head/foot
\IfFileExists{footnotehyper.sty}{\usepackage{footnotehyper}}{\usepackage{footnote}}
\makesavenoteenv{longtable}
\setlength{\emergencystretch}{3em} % prevent overfull lines
\providecommand{\tightlist}{%
  \setlength{\itemsep}{0pt}\setlength{\parskip}{0pt}}

\usepackage{newunicodechar}
\newunicodechar{→}{\ensuremath{\rightarrow}}
\newunicodechar{↔}{\ensuremath{\leftrightarrow}}
\usepackage{fvextra}
\DefineVerbatimEnvironment{verbatim}{Verbatim}{breaklines=true,breakanywhere=true,fontsize=\small}
\usepackage{fancyhdr}
\usepackage{array}
\pagestyle{fancy}\fancyhf{}
\fancyhead[L]{\small MovieLOD · Bài làm hoàn chỉnh}
\fancyfoot[C]{\small\thepage}\setlength{\headheight}{15pt}
\renewcommand{\arraystretch}{1.14}\setlength{\extrarowheight}{1pt}
\AtBeginEnvironment{longtable}{\small}
\setlength{\emergencystretch}{3em}\setlength{\parskip}{4pt}
\usepackage{titling}\setlength{\droptitle}{-1.5cm}
\pretitle{\begin{center}\Large\bfseries\color{teal}}
\posttitle{\par\end{center}\vspace{-0.5em}}
\predate{\begin{center}\small}\postdate{\par\end{center}\vspace{-1em}}
\usepackage{bookmark}
\IfFileExists{xurl.sty}{\usepackage{xurl}}{} % add URL line breaks if available
\urlstyle{same}
% fallback for those not using the hyperref driver hyperxmp:
\makeatletter
\@ifundefined{xmpquote}{\newcommand{\xmpquote}[1]{#1}}{}
\makeatother
\hypersetup{
  pdftitle={MovieLOD: hướng dẫn thao tác chi tiết bản 2.0},
  colorlinks=true,
  linkcolor={teal},
  filecolor={Maroon},
  citecolor={Blue},
  urlcolor={teal},
  pdfcreator={LaTeX via pandoc}}

\title{MovieLOD: hướng dẫn thao tác chi tiết bản 2.0}
\author{}
\date{MovieLOD 2.0 · Đối chiếu 08/10/2026}

\begin{document}
\maketitle

\subsection{1. Phiên bản, thứ tự đọc và kết quả cần
biết}\label{phiuxean-bux1ea3n-thux1ee9-tux1ef1-ux111ux1ecdc-vuxe0-kux1ebft-quux1ea3-cux1ea7n-biux1ebft}

Tài liệu cho \textbf{ontology 2.0}, cập nhật 08/10/2026. Giao diện tiếng
Anh, hướng dẫn tiếng Việt. Đọc
Huong\_\allowbreak{}da\allowbreak{}n\_\allowbreak{}\allowbreak{}A\_\allowbreak{}\allowbreak{}Z
trước; tra bảng lớp ở
Ontology\allowbreak{}\_\allowbreak{}\allowbreak{}Redesig\allowbreak{}n;
tập nói theo
Script\_\allowbreak{}t\allowbreak{}huyet\_\allowbreak{}tr\allowbreak{}inh;
quay theo
Kich\_\allowbreak{}ban\allowbreak{}\_\allowbreak{}video.\allowbreak{}

\textbf{Mốc local:} 42 lớp, 30 phim, 851 người, 1.010 đóng góp, 45 công
ty, 672 thực thể giải, 19.339 triple dữ liệu, 518 triple lược đồ, 1.727
sameAs. Có 24 file SPARQL và 14 mẫu trực tiếp trên Web. Số lượng có thể
thay đổi nếu tải lại nguồn; phải đo lại thay vì dùng mốc này cho phiên
bản mới.

\textbf{Các lỗi đã sửa:} bổ sung 43 file còn thiếu, hiện đủ 76/76 và
hash khớp; công bố bản 2.0 với 19.339 triple và 42 lớp, graph public
đẳng cấu local. 14 test pass vẫn cần đi kèm kiểm tra nguồn. Biên bản mới
ở
evidence\allowbreak{}/\allowbreak{}review\_\allowbreak{}2026-\allowbreak{}10-\allowbreak{}08.\allowbreak{}json
và
publicat\allowbreak{}ion\_\allowbreak{}chec\allowbreak{}ks.\allowbreak{}json.\allowbreak{}
Video\_\allowbreak{}de\allowbreak{}mo.\allowbreak{}mp4 đã được thay bằng
demo tương tác bản 2.0 với lời tiếng Việt tổng hợp.

\subsection{2. Chuẩn bị môi trường và chạy
ngay}\label{chuux1ea9n-bux1ecb-muxf4i-trux1b0ux1eddng-vuxe0-chux1ea1y-ngay}

Mở terminal tại thư mục chứa README.md và config.json:

\begin{verbatim}
pwd
python3 --version
python3 -m venv .venv
.venv/bin/python -m pip install -r requirements.txt
.venv/bin/python src/server.py
\end{verbatim}

Mở
\textbf{http:\allowbreak{}/\allowbreak{}/\allowbreak{}1\allowbreak{}27.\allowbreak{}0.\allowbreak{}0.\allowbreak{}1\allowbreak{}:\allowbreak{}8000}.
Dừng bằng Ctrl+C. Có thể dùng start.\allowbreak{}co\allowbreak{}mmand
trên macOS. Nếu cổng bị chiếm, dùng
\texttt{-\allowbreak{}-\allowbreak{}port 8\allowbreak{}001}; thay cổng
trong mọi URL/curl. Không tự dừng server của người khác. Server nạp dữ
liệu lúc khởi động; \textbf{sau build phải khởi động lại} để endpoint
không giữ graph cũ trong RAM.

Dữ liệu chuẩn hóa đã có; không cần chạy collect hay validate để mở ứng
dụng. Khi ghi video hãy dùng một server vừa khởi động trên cổng trống.

\subsection{3. Kiểm tra giao diện lần
đầu}\label{kiux1ec3m-tra-giao-diux1ec7n-lux1ea7n-ux111ux1ea7u}

\begin{enumerate}
\def\labelenumi{\arabic{enumi}.}
\tightlist
\item
  Chọn \textbf{Inception: year, runtime and director}; bấm \textbf{Run
  query} nếu cần. Thấy Inception, 2010, 148, Christopher Nolan.
\item
  Chọn \textbf{Films directed by Christopher Nolan}: 8 dòng trong
  dataset.
\item
  Chọn \textbf{Who contributed to Inception, and in which role?}: 25
  dòng, các vai trò
  Actor/\allowbreak{}\allowbreak{}Di\allowbreak{}rector/\allowbreak{}\allowbreak{}W\allowbreak{}riter/\allowbreak{}\allowbreak{}Pr\allowbreak{}oducer.\allowbreak{}
  Hai vai trò cùng người có các Contribution riêng.
\item
  Chọn \textbf{Production companies of Inception}: 4 dòng.
  \textbf{Awards received by The Godfather}: 7 dòng.
\item
  Chọn \textbf{Is Inception linked to its Wikidata entity?}: True.
\item
  Bấm \textbf{Download results} để tải kết quả. Dùng Find a film để tìm
  Inception, mở trang thực thể, nguồn và link Download RDF (Turtle).
\end{enumerate}

Các số option không cố định khi thêm câu hỏi; dùng \textbf{nhãn}. Query
editor có thể sửa SPARQL. Gõ
\texttt{not \allowbreak{}S\allowbreak{}P\allowbreak{}A\allowbreak{}R\allowbreak{}\allowbreak{}Q\allowbreak{}L},
bấm Run query để xem báo lỗi; chọn lại câu hỏi mẫu để phục hồi.

\subsection{4. Đọc dữ liệu nguồn và hiểu
YC2}\label{ux111ux1ecdc-dux1eef-liux1ec7u-nguux1ed3n-vuxe0-hiux1ec3u-yc2}

Mở src/\allowbreak{}coll\allowbreak{}ect.\allowbreak{}py:\allowbreak{}
phim được tìm bằng sitelink Wikipedia chính xác. Wikidata cung cấp dữ
liệu về đạo diễn, diễn viên, biên kịch, nhà sản xuất, thể loại, quốc
gia, ngôn ngữ, năm, thời lượng, giải và công ty. DBpedia chỉ nối khi
đúng chủ thể Film.

Mở
data/\allowbreak{}raw\allowbreak{}/\allowbreak{}snapsho\allowbreak{}ts.\allowbreak{}json
để xem
U\allowbreak{}R\allowbreak{}L/\allowbreak{}prov\allowbreak{}ider/\allowbreak{}ret\allowbreak{}rieved\_\allowbreak{}a\allowbreak{}t/\allowbreak{}http\_\allowbreak{}s\allowbreak{}tatus/\allowbreak{}sh\allowbreak{}a256/\allowbreak{}pat\allowbreak{}h.\allowbreak{}
Chọn một row có path đang tồn tại để mở phản hồi gốc. Hash là dấu vân
tay byte; cùng hash chứng minh nội dung khớp với mốc lưu, không chứng
minh độ đúng ngoài đời.

\textbf{Nguồn đã hoàn tất:} danh mục có 76 row, 76 file có và 76 hash
khớp. validate.py hoàn tất và có biên bản mới. Một phản hồi đổi nội dung
được ghi hash/thời điểm mới; danh mục trước khôi phục được giữ để truy
lịch sử.

Khôi phục từ bản sao lưu đúng snapshot nếu có. Nếu phải tải lại, giữ mốc
thời điểm/hash mới, build và kiểm tra lại; dữ liệu Internet có thể thay
đổi. collect.py dùng cache đủ và khớp hash; khi thiếu sẽ tải,
\texttt{-\allowbreak{}-\allowbreak{}refres\allowbreak{}h} buộc tải lại.
\textbf{Chạy thu thập có thể thay đổi bộ dữ liệu}, hãy giữ bản sao của
data trước khi chủ động cập nhật.

\begin{verbatim}
.venv/bin/python src/collect.py
# Refresh all source responses:
.venv/bin/python src/collect.py --refresh
\end{verbatim}

\subsection{5. Sinh mô hình/RDF và đồng bộ Web
local}\label{sinh-muxf4-huxecnhrdf-vuxe0-ux111ux1ed3ng-bux1ed9-web-local}

\begin{verbatim}
.venv/bin/python src/build.py
.venv/bin/python src/prepare_web.py
\end{verbatim}

Build dùng dữ liệu chuẩn hóa hiện có và sinh RDF, OWL, số liệu, trang
thực thể và dữ liệu cho Web. Prepare\_web cập nhật truy vấn mẫu/giấy
phép/PDF trong web/dist. Không cần tải lại nguồn để build bản đã chuẩn
hóa. Sau khi build, khởi động lại server.

{\def\LTcaptype{none} % do not increment counter
\begin{longtable}[]{@{}
  >{\raggedright\arraybackslash}p{(\linewidth - 2\tabcolsep) * \real{0.5000}}
  >{\raggedright\arraybackslash}p{(\linewidth - 2\tabcolsep) * \real{0.5000}}@{}}
\toprule\noalign{}
\begin{minipage}[b]{\linewidth}\raggedright
File
\end{minipage} & \begin{minipage}[b]{\linewidth}\raggedright
Dùng để làm gì
\end{minipage} \\
\midrule\noalign{}
\endhead
\bottomrule\noalign{}
\endlastfoot
ontology\allowbreak{}/\allowbreak{}\allowbreak{}Movie\_\allowbreak{}\allowbreak{}O\allowbreak{}ntology.\allowbreak{}owl
& Chỉ mô hình; mở trong Protégé để xem lớp/ràng buộc \\
ontology\allowbreak{}/\allowbreak{}\allowbreak{}Movie\_\allowbreak{}\allowbreak{}K\allowbreak{}nowledge\allowbreak{}\_\allowbreak{}\allowbreak{}Graph.\allowbreak{}o\allowbreak{}wl
& Mô hình + dữ liệu; xem cá thể \\
ontology\allowbreak{}/\allowbreak{}movie.\allowbreak{}t\allowbreak{}tl &
Lược đồ Turtle; nguồn đối chiếu 42 lớp \\
data/\allowbreak{}pro\allowbreak{}cessed/\allowbreak{}m\allowbreak{}ovies.\allowbreak{}tt\allowbreak{}l
& Graph khai báo; endpoint truy vấn file này \\
data/\allowbreak{}pro\allowbreak{}cessed/\allowbreak{}m\allowbreak{}ovies.\allowbreak{}js\allowbreak{}onld
& Cùng dữ liệu dưới dạng JSON-LD \\
data/\allowbreak{}pro\allowbreak{}cessed/\allowbreak{}i\allowbreak{}nferred\_\allowbreak{}classes.\allowbreak{}ttl
& Các kiểu phân loại bổ sung, tách khỏi graph gốc \\
web/\allowbreak{}dist\allowbreak{}/\allowbreak{}data/\allowbreak{} & Dữ
liệu cấp cho engine trình duyệt \\
evidence\allowbreak{}/\allowbreak{}statist\allowbreak{}ics.\allowbreak{}json\allowbreak{}
& Số liệu mới sau build \\
\end{longtable}
}

Build không khắc phục các snapshot gốc đã mất. ``Dựng được từ dữ liệu
chuẩn hóa'' khác ``tái lập được từ mọi phản hồi gốc''.

\subsection{6. Xem ontology trong
Protégé}\label{xem-ontology-trong-protuxe9guxe9}

\begin{enumerate}
\def\labelenumi{\arabic{enumi}.}
\tightlist
\item
  File → Open →
  ontology\allowbreak{}/\allowbreak{}\allowbreak{}Movie\_\allowbreak{}\allowbreak{}O\allowbreak{}ntology.\allowbreak{}owl.\allowbreak{}
  Trong Entities → Classes, chọn Film, Person, Country và kiểm tra IRI
  DBpedia.
\item
  Mở
  Creative\allowbreak{}\allowbreak{}Work/\allowbreak{}\allowbreak{}Age\allowbreak{}nt/\allowbreak{}\allowbreak{}Contr\allowbreak{}ibution/\allowbreak{}\allowbreak{}Genre/\allowbreak{}\allowbreak{}Aw\allowbreak{}ard
  để xem nhánh. Số lớp 42 đếm owl:Class IRI tự khai báo; Protégé có thể
  hiển thị thêm lớp tham chiếu ngoài và biểu thức vô danh.
\item
  Object properties →
  contribu\allowbreak{}tion\allowbreak{}By/\allowbreak{}c\allowbreak{}ontribut\allowbreak{}ion\allowbreak{}To/\allowbreak{}ha\allowbreak{}s\allowbreak{}Role:\allowbreak{}
  kiểm tra domain, range, functional; lớp Contribution có restriction
  đúng một người/phim/vai trò.
\item
  Chọn
  has\allowbreak{}Contr\allowbreak{}ibution/\allowbreak{}contribu\allowbreak{}tion\allowbreak{}Of
  để xem inverse; director có inverse directed.
\item
  Mở
  Movie\_\allowbreak{}\allowbreak{}Kn\allowbreak{}owledge\_\allowbreak{}\allowbreak{}Graph.\allowbreak{}ow\allowbreak{}l,
  Individuals chọn film-Q25188 (Inception),
  person-\allowbreak{}\allowbreak{}Q\allowbreak{}25191 (Nolan), một
  Contribution của Nolan. Xem vai trò, phim và nguồn.
\item
  Chọn Filmmake\allowbreak{}r/\allowbreak{}\allowbreak{}Actor để xem
  equivale\allowbreak{}nt\allowbreak{}Class.\allowbreak{} Các kiểu không
  gán sẵn cần reasoning mới xuất hiện.
\end{enumerate}

\textbf{Giới hạn:}
Hermi\allowbreak{}T/\allowbreak{}\allowbreak{}P\allowbreak{}ellet đã
chạy trên
Movie\_\allowbreak{}\allowbreak{}Kn\allowbreak{}owledge\_\allowbreak{}\allowbreak{}Graph.\allowbreak{}ow\allowbreak{}l
và xác nhận nhất quán. Timestamp trong bản RDF chuẩn hóa dùng mili giây;
bản nguồn thô giữ thời điểm đầy đủ. OWL không mặc định hai IRI là cá thể
khác nhau; không khẳng định ngưỡng cardinality bằng số IRI là chứng minh
DL.
Domain/\allowbreak{}r\allowbreak{}ange/\allowbreak{}car\allowbreak{}dinality\allowbreak{}
không tự thay kiểm tra biểu mẫu bằng Python.

\subsection{7. Kiểm tra và tạo phân
loại}\label{kiux1ec3m-tra-vuxe0-tux1ea1o-phuxe2n-loux1ea1i}

\begin{verbatim}
.venv/bin/python src/validate.py
.venv/bin/python src/reason.py
.venv/bin/python -m pytest -q
\end{verbatim}

Validate kiểm tra các trường cơ bản, hash phản hồi nguồn và chạy file
SPARQL. Lệnh hiện chạy hoàn tất: data\_checks\_passed=true,
source\_hashes\_match=true, snapshots\_checked=76 và không có file
thiếu. Khi tải lại phải đo mới số lượng/hash thay vì mặc định mốc cũ.

Reason thực hiện OWL RL cho 12 lớp, COUNT DISTINCT cho 2 lớp
Multi\allowbreak{}Gen\allowbreak{}re\allowbreak{}Film/\allowbreak{}\allowbreak{}F\allowbreak{}ilm\allowbreak{}Studi\allowbreak{}o.\allowbreak{}
Kết quả ở
ontology\allowbreak{}\_\allowbreak{}reasoni\allowbreak{}ng.\allowbreak{}json
và
inferred\allowbreak{}\_\allowbreak{}classes\allowbreak{}.\allowbreak{}ttl.\allowbreak{}
Không chạy trong thời gian demo nếu mất nhiều thời gian; chạy trước để
tạo file. Không có lỗi được phát hiện bởi owlrl không đồng nghĩa chứng
minh OWL DL nhất quán.

Test hiện \textbf{14 passed} (xem thời gian thực tế trong
evidence\allowbreak{}/\allowbreak{}tests.\allowbreak{}t\allowbreak{}xt).\allowbreak{}
Có kiểm tra endpoint, kiểu DBpedia, chuẩn hóa Inception, role, kiểm tra
dữ liệu thiếu, OWL RL và các định nghĩa lớp. Chạy test là một phần xác
minh; kiểm tra 76 nguồn được thực hiện riêng và đã đạt.

\subsection{8. Terminal và endpoint}\label{terminal-vuxe0-endpoint}

\textbf{Terminal đọc graph gốc:}

\begin{verbatim}
.venv/bin/python src/query.py queries/02_inception.rq
.venv/bin/python src/query.py queries/04_credits.rq
.venv/bin/python src/query.py queries/08_ask.rq
\end{verbatim}

\textbf{Truy vấn có schema/phân loại:}

\begin{verbatim}
.venv/bin/python src/query.py queries/13_fiction_genre_films_via_hierarchy.rq --reasoned
.venv/bin/python src/query.py queries/18_inferred_filmmakers.rq --reasoned
.venv/bin/python src/query.py queries/24_nolan_asserted_types_only.rq
.venv/bin/python src/query.py queries/24_nolan_asserted_types_only.rq --reasoned
\end{verbatim}

Câu 13 trả 29 phim khi có schema. Câu 18 hiển thị 20 dòng do LIMIT dù
tổng có 89 Filmmaker. Câu 24 gốc có 1 kiểu; có phân loại có 3 kiểu.
--reasoned nạp file phân loại đã có, không chạy suy luận ngay.

\textbf{Endpoint --- server đang chạy ở terminal A:}

\begin{verbatim}
curl -G http://127.0.0.1:8000/sparql \
  --data-urlencode query@queries/02_inception.rq
curl -X POST http://127.0.0.1:8000/sparql \
  -H 'Content-Type: application/sparql-query' \
  --data-binary @queries/02_inception.rq
curl -L -H 'Accept: text/turtle' \
  http://127.0.0.1:8000/resource/film-Q25188
\end{verbatim}

GET/POST SELECT/ASK trả JSON.
C\allowbreak{}O\allowbreak{}N\allowbreak{}S\allowbreak{}T\allowbreak{}R\allowbreak{}U\allowbreak{}C\allowbreak{}\allowbreak{}T/\allowbreak{}\allowbreak{}D\allowbreak{}E\allowbreak{}S\allowbreak{}C\allowbreak{}R\allowbreak{}I\allowbreak{}\allowbreak{}B\allowbreak{}E
trả Turtle. Accept Turtle tại resource cục bộ trả 303; -L theo chuyển
hướng. Hosted dùng Comunica trong trình duyệt và file RDF, không phải
endpoint Flask public. SERVICE/FROM tải ngoài và SPARQL Update bị chặn.

\subsection{9. Ảnh minh chứng và kiểm tra trình
duyệt}\label{ux1ea3nh-minh-chux1ee9ng-vuxe0-kiux1ec3m-tra-truxecnh-duyux1ec7t}

Script bổ sung cần Playwright và Google Chrome đã cài, hoặc Chromium của
Playwright. Script không bắt buộc để mở ứng dụng.

\begin{verbatim}
.venv/bin/python -m pip install playwright
# If a suitable Chrome is not installed:
.venv/bin/python -m playwright install chromium
.venv/bin/python src/browser_check.py --port 8000
\end{verbatim}

Server phải chạy đúng cổng. Có thể dùng --port 8001/8002 khi kiểm tra.
Script ghi 4 ảnh
(app/\allowbreak{}\allowbreak{}Nol\allowbreak{}an/\allowbreak{}resou\allowbreak{}rce/\allowbreak{}mobi\allowbreak{}le)
và
browser\_\allowbreak{}checks.\allowbreak{}j\allowbreak{}son.\allowbreak{}
Minh chứng: SELECT endpoint, Comunica SELECT/GROUP
B\allowbreak{}Y/\allowbreak{}\allowbreak{}A\allowbreak{}S\allowbreak{}K/\allowbreak{}\allowbreak{}C\allowbreak{}\allowbreak{}O\allowbreak{}N\allowbreak{}S\allowbreak{}T\allowbreak{}R\allowbreak{}U\allowbreak{}C\allowbreak{}T\allowbreak{},
thông báo lỗi, trang IRI, không tràn màn hình mobile, không có exception
trình duyệt.

Ảnh trong slide phải lấy \textbf{sau khi build và khởi động lại server}.
Nếu chỉ dữ liệu endpoint mới nhưng file web cũ, KPI và trang thực thể sẽ
lệch; build + prepare\_web đồng bộ file local rồi chụp lại. Việc này
chưa xuất bản lên Internet.

\subsection{10. Sửa tài liệu, tạo slide và
PDF}\label{sux1eeda-tuxe0i-liux1ec7u-tux1ea1o-slide-vuxe0-pdf}

Nguồn nội dung tài liệu là các file Markdown trong docs. Lời slide nằm
trong
src/\allowbreak{}pres\allowbreak{}entation\allowbreak{}\_\allowbreak{}content\allowbreak{}.\allowbreak{}py;
sơ đồ và bố cục trong
src/\allowbreak{}make\allowbreak{}\_\allowbreak{}slides\_\allowbreak{}video.\allowbreak{}py\allowbreak{}.\allowbreak{}
Font chung \textbf{Be Vietnam Pro}. Các công cụ PDF cần Pandoc, Tectonic
và font; máy hiện tại đã có.

\begin{verbatim}
.venv/bin/python src/make_slides_video.py --slides-only
.venv/bin/python src/make_docs.py
.venv/bin/python src/prepare_web.py
\end{verbatim}

make\_docs chỉ render Markdown hiện có, không tự ghi đè nội dung đã sửa.
Có thể render riêng:
\texttt{.\allowbreak{}venv/\allowbreak{}bi\allowbreak{}n/\allowbreak{}python\allowbreak{} src/\allowbreak{}mak\allowbreak{}e\_\allowbreak{}docs.\allowbreak{}p\allowbreak{}y docs/\allowbreak{}\allowbreak{}B\allowbreak{}ao\_\allowbreak{}cao.\allowbreak{}m\allowbreak{}d}.
Huong\_\allowbreak{}da\allowbreak{}n\_\allowbreak{}thao\_\allowbreak{}t\allowbreak{}ac\_\allowbreak{}chi\_\allowbreak{}t\allowbreak{}iet
cũng có thể render bằng
make\_\allowbreak{}man\allowbreak{}ual.\allowbreak{}py.\allowbreak{} PDF
slide hiện có 24 trang; PowerPoint có editable text/diagram và Speaker
Notes. PDF slide hiện là bản vector xuất từ PPTX qua LibreOffice; thiếu
LibreOffice thì trình tạo dùng bản raster dự phòng.

Nếu máy người đọc không có Be Vietnam Pro, cài font để PPTX giữ đúng bố
cục hoặc dùng PDF. Điền nhóm/thành viên/lớp trên bìa. Không cần thêm ảnh
stock; slide có sơ đồ mô hình và ảnh ứng dụng thật. Có thể thêm ảnh
Protégé thật theo cảnh YC1 trong kịch bản nếu muốn.

\subsection{11. Quay video 3--5 phút}\label{quay-video-35-phuxfat}

Dùng
Kich\_\allowbreak{}ban\allowbreak{}\_\allowbreak{}video.\allowbreak{}m\allowbreak{}d/\allowbreak{}pdf:\allowbreak{}
timeline 0:00--4:50, lời đọc từng cảnh, lệnh thật và kết quả mong đợi.
Script\_\allowbreak{}t\allowbreak{}huyet\_\allowbreak{}tr\allowbreak{}inh.\allowbreak{}md/\allowbreak{}p\allowbreak{}df
là lời nói 24 slide (\textasciitilde18--22 phút), không dùng nguyên để
quay video 3--5 phút.

Quay thao tác Protégé/editor/Web/terminal. Giữ cảnh nguồn/RDF/giấy
phép/sameAs; demo không chỉ bấm hai câu hỏi.
Video\_\allowbreak{}de\allowbreak{}mo.\allowbreak{}mp4 hiện là bản demo
2.0 dài khoảng 4:50, ghi thao tác trình duyệt và kết quả lệnh thật, có
lời tiếng Việt tổng hợp. Xem lại trước khi nộp; có thể thay bằng giọng
của thành viên.

\begin{verbatim}
ffprobe -v error -show_entries format=duration \
  -of default=noprint_wrappers=1:nokey=1 docs/Video_demo_2_0.mp4
\end{verbatim}

\subsection{12. Công bố và đối chiếu bản
public}\label{cuxf4ng-bux1ed1-vuxe0-ux111ux1ed1i-chiux1ebfu-bux1ea3n-public}

Lần kiểm tra mới 08/10: trang chủ, dataset,
Turtle/\allowbreak{}\allowbreak{}J\allowbreak{}\allowbreak{}S\allowbreak{}O\allowbreak{}N-\allowbreak{}\allowbreak{}L\allowbreak{}D,
ontology, IRI Inception, mô tả RDF và giấy phép đều trả 200 không đăng
nhập. RDF public có 19.339 triple, ontology 42 lớp; cả 4 graph RDF kiểm
tra đều đẳng cấu với local. Kết luận dùng kết quả HTTP và graph, không
chỉ trường audience\allowbreak{}/\allowbreak{}status trong JSON.

Sau khi nhóm xuất bản bản mới, kiểm tra từ phiên không đăng nhập: trang
dataset,
Turtle/\allowbreak{}\allowbreak{}J\allowbreak{}\allowbreak{}S\allowbreak{}O\allowbreak{}N-\allowbreak{}\allowbreak{}L\allowbreak{}D,
ontology, IRI và giấy phép. Tải RDF public về và so sánh \textbf{graph
đẳng cấu}, không so byte vì thứ tự Turtle có thể khác. Lưu thời điểm,
status, URL và kết quả thật trong evidence. Hosted có file tĩnh nên
không mặc định có hành vi 303/content negotiation như Flask local.

Đổi tên miền cần dùng
src/\allowbreak{}reba\allowbreak{}se.\allowbreak{}py rồi
build/\allowbreak{}pr\allowbreak{}epare\_\allowbreak{}we\allowbreak{}b;
IRI trong dữ liệu/truy vấn/giao diện phải thống nhất. Site đã được triển
khai lại qua quy trình hosting; trạng thái và ID phiên bản thật nằm
trong
publicat\allowbreak{}ion.\allowbreak{}json\allowbreak{}.\allowbreak{}
Khi sửa sau này, xuất bản và đối chiếu lại.

\subsection{13. Chấm và kiểm tra sản phẩm
nộp}\label{chux1ea5m-vuxe0-kiux1ec3m-tra-sux1ea3n-phux1ea9m-nux1ed9p}

Theo thang tự đề xuất chia đều 2 điểm mỗi YC: cả 5 YC đều đạt 2 điểm;
tổng \textbf{10/10}. Đề gốc không có trọng số; đọc CHAM\_DIEM.md để xem
tiêu chí con và minh chứng. Đây không phải điểm chính thức của giảng
viên.

Báo cáo tối đa 15 trang; slide 24 trang theo yêu cầu nhóm; video mới
180--300 giây. Có mã nguồn, RDF, OWL, dữ liệu gốc đủ, giấy phép, scripts
và evidence. Không đóng gói
.\allowbreak{}venv/\allowbreak{}\allowbreak{}Gi\allowbreak{}t/\allowbreak{}cache;\allowbreak{}
kiểm tra file nguồn có thực sự trong ZIP vì mọi phản hồi gốc phải nằm
trong bộ nộp; quy tắc gitignore làm mất nguồn đã được bỏ.
\texttt{src/\allowbreak{}pack\allowbreak{}age.\allowbreak{}py} là tiện
ích đóng gói, không thay các thiếu sót thành đạt.

\subsection{14. Lỗi thường gặp}\label{lux1ed7i-thux1b0ux1eddng-gux1eb7p}

{\def\LTcaptype{none} % do not increment counter
\begin{longtable}[]{@{}
  >{\raggedright\arraybackslash}p{(\linewidth - 2\tabcolsep) * \real{0.5000}}
  >{\raggedright\arraybackslash}p{(\linewidth - 2\tabcolsep) * \real{0.5000}}@{}}
\toprule\noalign{}
\begin{minipage}[b]{\linewidth}\raggedright
Hiện tượng
\end{minipage} & \begin{minipage}[b]{\linewidth}\raggedright
Cách kiểm tra và xử lý
\end{minipage} \\
\midrule\noalign{}
\endhead
\bottomrule\noalign{}
\endlastfoot
Port already in use & Dùng --port cổng trống và sửa URL/curl; không dừng
tiến trình khác \\
Contribution trả 0 & Server giữ graph cũ;
build/\allowbreak{}pr\allowbreak{}epare\_\allowbreak{}we\allowbreak{}b
và khởi động server mới \\
KPI/ảnh còn 15 lớp hoặc 964 links & File web cũ; build lại để đồng bộ
statisti\allowbreak{}cs.\allowbreak{}json và RDF \\
validate
File\allowbreak{}Not\allowbreak{}F\allowbreak{}ound\allowbreak{}Erro\allowbreak{}r
& Thiếu snapshot; bổ sung bản gốc hoặc thu thập lại có mốc mới \\
Filmmake\allowbreak{}r/\allowbreak{}\allowbreak{}Actor trả 0 trên
endpoint & Chạy terminal --reasoned; endpoint mặc định chỉ graph gốc \\
Câu hỏi chỉ 20 người & LIMIT 20; tổng
Actor/\allowbreak{}\allowbreak{}Fi\allowbreak{}lmmaker nằm ở
ontology\allowbreak{}\_\allowbreak{}reasoni\allowbreak{}ng.\allowbreak{}json \\
Documentary ASK false & Mẫu chưa có phim tài liệu; không phải kết luận
ngoài dataset \\
PDF lỗi công cụ/font & Kiểm tra
Pandoc/\allowbreak{}\allowbreak{}T\allowbreak{}ectonic/\allowbreak{}\allowbreak{}Be
Vietnam Pro, đọc build\_*.log \\
PPTX chữ lệch máy khác & Cài Be Vietnam Pro hoặc trình chiếu PDF \\
Hosted có 200 nhưng số liệu cũ & Xuất bản đúng bản mới rồi đối chiếu
graph \\
\end{longtable}
}

\subsection{15. Bổ sung ảnh Protégé vào bộ 24
slide}\label{bux1ed5-sung-ux1ea3nh-protuxe9guxe9-vuxe0o-bux1ed9-24-slide}

Có 11 khung P00--P10, có hướng dẫn ngay trên slide và Speaker Notes. Đọc
Checklis\allowbreak{}t\_\allowbreak{}anh\_\allowbreak{}\allowbreak{}Pr\allowbreak{}otege.\allowbreak{}md\allowbreak{}/\allowbreak{}pdf
để chọn đúng file OWL, lớp/thuộc tính/cá thể và nội dung cần thấy. Khung
chờ chưa phải ảnh thật. Có thể chèn trực tiếp vào PPTX hoặc lưu ảnh đúng
stem ở evidence\allowbreak{}/\allowbreak{}protege\allowbreak{}
(\allowbreak{}P\allowbreak{}N\allowbreak{}G/\allowbreak{}\allowbreak{}J\allowbreak{}P\allowbreak{}G\allowbreak{}/\allowbreak{}\allowbreak{}J\allowbreak{}P\allowbreak{}E\allowbreak{}G),
chạy lại
make\_\allowbreak{}sli\allowbreak{}des\_\allowbreak{}vide\allowbreak{}o.\allowbreak{}py
--slides-only. Ảnh tự thay khung chờ khi tạo lại. Bản ngắn 13 trang vẫn
ở
Slide\_\allowbreak{}ng\allowbreak{}an\_\allowbreak{}13.\allowbreak{}pp\allowbreak{}tx/\allowbreak{}pdf.\allowbreak{}

\textbf{Cập nhật
timestam\allowbreak{}p/\allowbreak{}reason\allowbreak{}er:\allowbreak{}}
dùng duy nhất
Movie\_\allowbreak{}\allowbreak{}Kn\allowbreak{}owledge\_\allowbreak{}\allowbreak{}Graph.\allowbreak{}ow\allowbreak{}l
cho graph đầy đủ.
Hermi\allowbreak{}T/\allowbreak{}\allowbreak{}P\allowbreak{}ellet đã
chạy; xem
Ket\_\allowbreak{}qua\_\allowbreak{}reasoner\allowbreak{}.\allowbreak{}pdf.\allowbreak{}
Video\_\allowbreak{}de\allowbreak{}mo.\allowbreak{}mp4 giữ nguyên theo
yêu cầu nhóm; timestamp trong dữ liệu mới giảm đến mili giây, nội dung
phim/quan hệ giữ nguyên, đã kiểm tra ở
video\_\allowbreak{}da\allowbreak{}taset\_\allowbreak{}co\allowbreak{}mpatibil\allowbreak{}ity.\allowbreak{}json\allowbreak{}.\allowbreak{}
Video chưa thể hiện kết quả reasoner mới và vẫn nhắc bộ slide ngắn 13
trang.

\end{document}
